\chapter{LANDASAN TEORI}


Tuliskan di sini \textbf{Pedoman Penulisan Landasan Teori:} \\
Uraikan teori-teori ilmiah dasar dan terapan yang mendukung objek analisis aplikasi mobile Anda:
\begin{itemize}
    \item \textbf{Mobile Computing}: Konsep dasar komputasi mobile, arsitektur client-server mobile, dan karakteristik mobile network.
    \item \textbf{Platform Sistem Operasi}: Teori dasar mengenai sistem operasi mobile yang relevan (seperti Android SDK, arsitektur iOS, cocoapods, dll.).
    \item \textbf{Alat Bantu Perancangan}: Landasan teori mengenai HIPO (Hierarchy Input Process Output), Flowchart program, dan standar antarmuka pengguna (\textit{User Interface/User Experience}).
    \item \textbf{Framework/Bahasa Pemrograman}: Bahasa pemrograman pendukung (seperti Java, Kotlin, Swift, Dart/Flutter, atau React Native).
\end{itemize}

